\subsection{Counting the cycles and specializing to ten terms}\label{subsec:ten-power}\label{sec:finite-inputs}

\begin{proof}[Proof of Theorem~\ref{thm:main}]
Choose the spatial scales as in Lemma~\ref{lem:inductive-valleys}. Next
take $A$ sufficiently large for all the ordinary and incoming tests
of Lemma~\ref{lem:valley-rates} and for the two compact brackets and
two compact incoming tests of Lemma~\ref{lem:compact-seed}. Use
$P_{A,\lambda_A}$ at exact tail balance to choose the incoming heights.
Keep a small closed interval about $\lambda_A$ on which all $2+6L+j$
ordinary brackets persist. The tail difference is strictly monotone
there by Lemma~\ref{lem:tail-inward}(iii), so its endpoint signs are
opposite. Write $\omega=(-1)^j$.
If $\omega=-1$, the change $(\widetilde x,\widetilde y,\widetilde t)
=(-x,y,-t)$ replaces $P(x)$ by $P(-x)$. It exchanges the two
half-profiles and preserves the axis ordinates. This multiplies each difference of intercepts by
$\omega$ and makes the highest odd coefficient positive.
A common positive scaling of $x$ and $y$ makes the even leading
coefficient one. It scales every axis intercept by the same factor.

Read the incoming tests in decreasing physical height. Before
normalization the positive-height list runs from $\omega$ to $+1$.
Multiplying by $\omega$ gives endpoints $\omega^2=+1$ and $\omega$,
with $6L+j$ changes. Under the same normalization the two compact
negative seed heights have signs $-\omega,\omega$. The negative-height
valley list, read in increasing depth $H>0$, originally runs from $+1$
to $\omega$. Multiplying these endpoints by $\omega$ gives
$(\omega,\omega^2)=(\omega,+1)$ at the two ends, with $4L+j$ changes. Write $h_{\min}>0$ for the smallest
scaled positive incoming test height and $H_{\min}>0$ for the smallest
scaled valley depth. These tests and the compact seed ordinates were
fixed before choosing $A$. Hence for large $A$ the physical ordinates
before the common positive normalization satisfy
\[
 -A^{d_0}H_{\min}<h_{\mathrm{seed},2}<h_{\mathrm{seed},1}<0<A^{d_0}h_{\min}.
\]

Choose $s>0$ so small that every test retains its sign when
$G_+(h)-G_-(h)$ is replaced by $G_+(h)-G_-(h+s)$. Choose a physical height
$h=h_{\rm deep}$ below every preceding test so that $|G_+(h)-h|<s/4$
and $|G_-(h+s)-(h+s)|<s/4$. Such a height exists by
Lemma~\ref{lem:tail-inward}(iv), and its shifted comparison is less than $-s/2$. Proposition~\ref{prop:fold-closure}
transfers this entire ordered list to ordinary upper-axis launches in
one polynomial vector field. In increasing order of upper launch,
the four groups have endpoint signs
\[
 \underbrace{(+1,\ldots,\omega)}_{\text{positive heights}},\quad
 \underbrace{(-\omega,\omega)}_{\text{compact seed}},\quad
 \underbrace{(\omega,\ldots,+1)}_{\text{negative valley heights}},\quad
 \underbrace{(-1,+1)}_{\text{deep and exterior}} .
\]
At the seed the sequence $\omega,-\omega,\omega$ gives two changes
for either parity of $j$. Its last sign is the first valley sign.
At the far end $+1,-1,+1$ gives the other two changes. In total the
incoming list gives
\[
 (6L+j)+2+(4L+j)+2=10L+2j+4
\]
brackets. The ordinary tests give two compact brackets and $6L+j$
additional brackets. The fold itself contributes one bracket on each
side of the crest and two simple matched orbits. All lie in disjoint
launch intervals, after slightly separating any shared strict endpoints.
Lemma~\ref{lem:hyperbolicity} gives the simultaneous count
\[
 (2+6L+j)+(10L+2j+4)+4=16L+10+3j .
\]

There are three seed exponents, $j$ prefix exponents, four per valley,
and three from the terminal fold. The even exponents increase from $d_0$
through $d_\nu+2D,d_\nu+3D$ to $d_L+8$. Among the odd exponents, the
last prefix power is at most $d_0+5$, the first valley power is
$d_0+2j+D-1$, and each valley ends one below its new leading even power.
The final odd exponents are $d_L+1,d_L+3$. Thus all $4L+6+j$ coefficients
occupy distinct powers. Each is nonzero by construction, and the final
regularization varies an existing coefficient in its nonzero interval.
Hyperbolicity gives an open coefficient set on this support.
Substituting $k=4L+6+j$ proves the bound and the lower limit
in \eqref{eq:main-bound}. In particular, taking $L=1$ and $j=0$ in this
analytic construction gives at least 26 cycles on ten terms for every
sufficiently large odd $n_0$.
\end{proof}

\begin{corollary}[A fixed ten-term support in degree $67088$]\label{cor:ten-support}
There exists a nonempty open set $V\subset(\mathbb R\setminus\{0\})^{10}$
such that every coefficient vector in $V$ gives at least $26$ simultaneous
hyperbolic limit cycles on the support
\begin{equation}\label{eq:ten-support}
 \{33153,66305,66306,66563,66822,67079,67080,67081,67083,67088\}.
\end{equation}
\end{corollary}
\begin{proof}
Take $L=1$, $j=0$ and $n_0=257$. Then $m_0=33153$, $d_0=66306$
and $d_L=67080$. The role of this choice is to fix the ten exponents
in \eqref{eq:ten-support}. For a general number of blocks,
Theorem~\ref{thm:main} obtains the required local inequalities by taking
$n_0$ sufficiently large. Here we check them at $257$.
The interval computations below enclose the seed quantities, the
positive-height variations and the five negative-height integrals,
using $768$-bit real balls in Arb \cite{Johansson2017}.
These are inputs to Lemmas~\ref{lem:compact-seed}
and~\ref{lem:inductive-valleys}.

At the seed, bisection in $T=r^{m_0}$ isolates the minimum, the inner
crest and both upper-level inverses in Lemma~\ref{lem:compact-seed}.
Evaluating the squared-slope antiderivatives in $768$-bit Arb balls gives
\[
 0.81<\frac{H_b}{d_0}<0.82,\qquad 4<\frac{S_+(d_0)}{m_0^3}<5,\qquad
 16<\frac{S_-(d_0)}{m_0^3}<17,\qquad \frac{q_-''(v)}{v^2q_+''(1)}>10^5 .
\]
At the valley, take $\eta=(2/3)^{257}$ and $q=(3/4)^{128}$. By Descartes'
rule of signs, the cubic factor $257-259(3-\eta)z^2+520z^3$ of $E'(r)$
has at most two positive roots; bisection encloses the crest $z_p$, the
minimum $z_c$, and the critical heights $m=E(c)$ and $M=E(p)$ within
\[
 \begin{aligned}
  z_p &\in (0.99229,\,0.99230), & 1-z_c &\in (2.39,\,2.40)\cdot 10^{-44}, \\
  m &\in (5.55,\,5.56)\cdot 10^{-46}, & M &\in (2.42,\,2.43)\cdot 10^{-5}.
 \end{aligned}
\]
Because $7/8<z_p$, the profile $E$ increases strictly on $[1/2,7/8]$;
bisection of $z^{257}C_\eta(z)=m$ on that interval encloses the unique
pre-crest inverse $z_a\in(0.67022,\,0.67023)$. At the six outer nodes
$z\in\{2^{-13},1/8,1/2,3/2,16,2^{13}\}$ of Lemma~\ref{lem:valley-local}
and at the post-valley node $z_v=z_c+2^{-200}$ (where $m<E(z_v^{1/D})<M$),
$768$-bit ball evaluation of the exact polynomial expressions for $V_o$
and $V_i$ gives the following enclosures (at $z_v$, both $V_o=P_1$ and
$V_i=P_1+U(z_v^{1/D})-U(c)$ lie in the listed negative interval):
\[
\begin{array}{@{}ccc@{}}
\toprule
 z & \text{Certified enclosure} & \text{Sign}\\
\midrule
 2^{-13} & (1.76,1.77)\cdot 10^{-2012} & +\\
 1/8 & (-5.11,-5.10)\cdot 10^{-464} & -\\
 1/2 & (2.07,2.08)\cdot 10^{-154} & +\\
 z_c+2^{-200} & (-3.71,-3.70)\cdot 10^{-57} & -\\
 3/2 & (5.71,5.72)\cdot 10^{92} & +\\
 16 & (-1.22,-1.21)\cdot 10^{626} & -\\
 2^{13} & (9.84,9.85)\cdot 10^{2033} & +\\
\bottomrule
\end{array}
\]
This certifies both seven-sign positive lists. For the five depths $H=w^{257}$ we enclose the normalized integral
$nHDJ(H)$ of \eqref{eq:J-in-z}, with $\nu=1/D$. Write its odd factor as
$B_q=(1-z)(8-z-z^2)/8-qz$. On $0<z\le1/8$,
$0<1-(9/8+q)z\le B_q<2$, $z^\nu\le1$, and the denominator is at least
$H^2$. The absolute normalized integrand is bounded by $2nz^{n-1}/H$.
On $z\ge8$, $C_\eta=1-(3-\eta)z^2+2z^3\ge(2-3/8)z^3>z^3$.
Since $z^3/8\ge8z$ and $q<1$, $0<B_q=1-(9/8+q)z+z^3/8\le z^3$.
Here the absolute integrand is at most $nHz^{-n+\nu-4}\le nHz^{-n-1}$.
Integration bounds the sum of the two absolute tails by $8^{-n}(2/H+H)$.

On each of $[1/8,1/4]$, $[1/4,1/2]$, $[3/2,2]$, $[2,4]$ and $[4,8]$
use $1024$ equal cells. On each side of $1$, subdivide the distances
$[2^{-l-1},2^{-l}]$ into $512$ equal cells for $1\le l\le99$ and
use one last cell from distance $2^{-100}$ to zero. On each of the
$106498$ cells the endpoint values enclose $(1-z)^2$ by monotonicity.
Use $C_\eta=(1-z)^2(1+2z)+\eta z^2$ and the factored $B_q$
above. Multiply the interval range of the normalized integrand by
the exact dyadic width, sum with outward rounding at $768$ bits in Arb,
and enlarge by the two-sided tail bound. Table~\ref{tab:main-depth-cert}
gives the resulting local enclosures. The strict enclosures allow us to choose the inner
dilation and balance the two limiting returns. For a sufficiently large
outer radius, Proposition~\ref{prop:fold-closure} then supplies the three
matching parameters. Lemma~\ref{lem:hyperbolicity} gives a regular value
of an existing odd coefficient. These existence arguments produce the
open coefficient set on \eqref{eq:ten-support} without specifying a
numerical vector of final coefficients. Table~\ref{tab:ordered-returns}
counts the $24$ sign brackets and two simple matched orbits. The
ten-term examples in Corollary~\ref{cor:short-lists} use the same
analytic closure with a degree-ten seed and the small-$\eta$ signs from
Proposition~\ref{prop:threshold}.
\end{proof}

\begin{table}[htbp]
  \centering
  \caption{Rational enclosures of $nHDJ(H)$ from \eqref{eq:J-in-z},
  with $n=257$, $D=258$ and $H=w^{257}$. Each enclosure incorporates analytic bounds
  for the tails $(0,1/8]$ and $[8,\infty)$ in the integration variable $z$.}\label{tab:main-depth-cert}
  \begin{tabular}{@{}ccc@{}}
    \toprule
    $w$ & Enclosure & Sign\\
    \midrule
    $1/2$ & $(7/10,\,6/5)$ & $+$\\
    $2/3$ & $(-4\cdot10^8,\,-3\cdot10^8)$ & $-$\\
    $7/8$ & $(12/5,\,3)$ & $+$\\
    $3/2$ & $(-29/100,\,-23/100)$ & $-$\\
    $4$ & $(1/25,\,7/100)$ & $+$\\
    \bottomrule
  \end{tabular}
\end{table}
